% MAS1004 Assignment 1: long report.
%
% You fill this in with your agent, in your repository. AGENTS.md tells the
% agent how. Red text in the PDF is what still has to be replaced.
%
% - Pictures go in figures/. Copy them from results/, which is not in git.
% - Outputs go in outputs/. run.py saves what it prints in
%   results/<tag>_output.txt, and PLAN.md saves the other outputs in results/
%   with `| tee results/...`. Copy them here. \outputfile{outputs/check.txt}
%   prints a file exactly as it is. Every number in this report has to come
%   from such an output.
%
% To make the PDF: compress this folder into a zip, and on Overleaf choose
% New Project, Upload Project. Set Menu, Compiler to XeLaTeX, and Recompile.

\documentclass[11pt]{article}
\usepackage{mas1004}

\reporttitle{Assignment 1: Long report}
\studentname{Your name}
\studentid{Your student ID}

\begin{document}
\makeheader

% ---------------------------------------------------------------------------
\section*{Problem 1. Categories and images}

\subsection*{Why these categories}
\fillin{Why these categories, and why you care about them.}

\subsection*{What I expected the model to use}
\fillin{What you thought, before training, the model would use to tell them
apart: shape, colour, texture, background, or something else.}

\subsection*{Images per category}
I downloaded 150 images for each category. For back\_lat\_spread I later
downloaded 59 more (Problem 4), so 209 in total, 659 for all four categories.
After cleaning, 361 were left: back\_double\_biceps 114, back\_lat\_spread 85,
side\_chest 86 and side\_triceps 76. The full count, as \texttt{clean.py count}
printed it, is in \texttt{outputs/clean\_count.txt}, shown under Problem 4.

% ---------------------------------------------------------------------------
\section*{Problem 2. Preparing an image}

\sidebyside{figures/prepare_original.png}{\fillin{the image as downloaded}}
           {figures/prepare_result.png}{\fillin{what prepare\_image made of it}}

\fillin{What was cut off, and an image where the crop cut off part of the
thing you care about.}

% ---------------------------------------------------------------------------
\section*{Problem 3. First training runs}

\begin{center}\footnotesize
\begin{tabular}{llllrrrr}
\toprule
tag & start & trained & epochs & lr & trainable parameters & train acc. & test acc. \\
\midrule
run & ImageNet & all layers & 10 & 0.0001 & 11,178,564 & 98.8\% & 83.3\% \\
scratch & random & all layers & 10 & 0.0001 & 11,178,564 & 83.3\% & 45.8\% \\
frozen & ImageNet & last layer & 10 & 0.001 & 2,052 & 86.9\% & 75.8\% \\
epochs20 & ImageNet & all layers & 20 & 0.0001 & 11,178,564 & 98.8\% & 85.8\% \\
clean & ImageNet & all layers & 10 & 0.0001 & 11,178,564 & 100.0\% & 97.2\% \\
\bottomrule
\end{tabular}
\end{center}

\picturefile{figures/run_curves.png}

\fillin{Did the loss go down? Did it keep going down, or flatten out?}

\fillin{The gap between train and test accuracy, and what it means.}

\begin{center}\picturefile[0.5\linewidth]{figures/run_confusion.png}\end{center}

\fillin{Which two categories it mixes up most, and whether that surprised you.}

\fillin{How far apart the test accuracies are, starting from ImageNet and from
random numbers, and why you think that is.}

\fillin{How close training only the last layer gets to training all of them.}

\subsection*{Output of the runs}
\outputfile{outputs/run_output.txt}
\outputfile{outputs/scratch_output.txt}
\outputfile{outputs/frozen_output.txt}
\outputfile{outputs/epochs20_output.txt}
\outputfile{outputs/clean_output.txt}

% ---------------------------------------------------------------------------
\section*{Problem 4. Cleaning}

\subsection*{My rule}
It is not an actual photograph of a person performing the pose — this excludes video thumbnails with text overlays, cartoon/diagram illustrations explaining a pose or muscle group, AI-generated images, and overview/collage images showing several different poses at once.

The person is wearing regular clothing rather than posing in competition or training attire.

More than one person is a main subject of the photo, making it unclear who the label refers to.

The pose shown is visibly a different pose than the category name — for example, a Front Lat Spread photo filed under Back Lat Spread, someone performing a gym exercise (lat pulldown, tricep pushdown) instead of holding the static pose, or a Back Double Biceps pose filed under Back Lat Spread.

The image is a near-duplicate of another image already kept in the same category (same photo, different crop/watermark/quality).

\textbf{Addition to my rule, made after seeing the suspect list:} For near-copies that appear across two different categories (the same photo filed under two different pose labels), I check which pose the photo actually shows. If it clearly matches one category, I keep it there and remove it from the other. If it is ambiguous which pose it shows, I remove it from both categories entirely, since keeping it in either would teach the model an unreliable label.

\textbf{Further additions, made while removing images after seeing the suspect list:} seven reasons I used do not match a line of the rule above as first written, so I add them here:
\begin{itemize}
  \item not a bodybuilding pose (stretch, not a competition pose) (2 images)
  \item wrong camera angle, pose not identifiable (6 images)
  \item not an actual photograph (low-quality screenshot of a social media post) (12 images)
  \item image combines a front and back view in one picture, does not match a single pose (1 image)
  \item image combines parts of two different people's bodies (1 image)
  \item image quality too poor to identify the pose (2 images)
  \item incomplete person visible in frame (1 image)
\end{itemize}

\subsection*{Three kinds of junk}
\sidebyside{figures/junk_thumbnail_back_lat_spread_0011.jpg}
           {1. A video thumbnail with text overlay (back\_lat\_spread/0011.jpg)}
           {figures/junk_wrong_pose_back_lat_spread_0022.jpg}
           {2. A wrong pose: back double biceps mislabeled as back lat spread
            (back\_lat\_spread/0022.jpg)}

\sidebyside{figures/junk_near_duplicate_side_chest_0046.jpg}
           {3. A near-duplicate, removed (side\_chest/0046.jpg)}
           {figures/junk_near_duplicate_kept_side_chest_0032.jpg}
           {the image kept from its group (side\_chest/0032.jpg), similarity 0.9803}

\subsection*{What I removed}
\outputfile{outputs/clean_count.txt}

The initial clean.py suspects run on the original 600 images found 139 near-copy pairs (this breakdown came from an ad-hoc command, not repository code): 35 same-category pairs at similarity $\geq$0.99, 78 same-category pairs at 0.95–0.99, 16 cross-category pairs at $\geq$0.99, and 10 cross-category pairs at 0.95–0.99.

For the 113 same-category pairs, I first dealt with the ones that also broke another rule (32 pairs had both images removed for other reasons, such as being video thumbnails or showing the wrong pose). Of the remaining 81 pairs, which formed 40 distinct groups, I kept exactly one image per group and removed the rest as near-duplicates (60 images removed in total). Two groups were an exception: in one, the two images (similarity 0.9775) turned out to be different photographs of different people, not copies, so both were kept; in another, two images were visually distinct frames from the same photo series (not directly flagged as similar to each other by the tool) and were both kept.

For the 26 cross-category pairs (the same photo filed under two different pose labels, often forming groups of 3–7 images across Side Chest and Side Triceps in particular), I applied my rule's addition: where the pose was clearly identifiable, I kept it in its correct category and removed it from the other(s) — this applied to 3 groups. The remaining 6 groups were removed entirely from all categories, though in every case for a specific reason (video thumbnail, stretch pose, regular clothing, or illustration) rather than because the pose itself was ambiguous.

Of the first 20 suspects in each class, I removed: back\_double\_biceps 9,
back\_lat\_spread 17, side\_chest 18, side\_triceps 17. These counts were
computed with an ad-hoc command, not with code in this repository, because
\texttt{clean.py count} does not break the removals down this way: it matched
the first 20 suspects of each class in
\texttt{results/cleaning/suspects/suspects.csv} against the files in
\texttt{data/removed/}.

After cleaning, back\_lat\_spread had 74 images left, far fewer than the
other three classes, so I downloaded more with the search phrase ``back lat
spread pose bodybuilding competition''. The download kept 150 images. 91 of
them were near copies (similarity above 0.95, the threshold \texttt{clean.py}
uses): 54 matched existing clean images, 35 matched previously removed
images, and 2 were near-duplicates of each other within the new download. I
left these 91 out. I added the
other 59 to \texttt{data/raw/back\_lat\_spread} and
\texttt{data/clean/back\_lat\_spread} as \texttt{0151.jpg} to \texttt{0209.jpg},
applied my rule to them, and removed 31 (30 at first, and \texttt{0180.jpg} later,
after it appeared among the clean model's most confident mistakes and I saw it
shows a front lat spread), which leaves 28. The 150, 91 (with its breakdown) and 59
were computed with an ad-hoc script, not with code in this repository, which
has no command for adding to one class; the 31 removed are in the
\texttt{clean.py count} output above. These 59 images were not in the suspect
list, so the first-20 counts above do not include them.

The suspect list caught most of the junk I had already noticed during an initial scan — mainly video thumbnails, illustrations, and the more obvious near-duplicates. It was not exhaustive, however. Going through every remaining image individually after clearing the suspects, I found substantially more junk it had not flagged: most notably images showing the wrong pose entirely (12 Front Lat Spread and 1 Back Double Biceps photo mislabeled as Back Lat Spread, for example), multiple-person photos, and low-quality social media screenshots. This second pass removed more images (115) than the first suspect-based pass (92). A separate near-copy sheet, built specifically to find same-category duplicate groups, then caught a further 60 duplicate images that neither the original suspect ranking nor my own manual scan had removed, since visually ``normal'' duplicate photos don't stand out the way genuinely odd or wrong images do. This suggests the suspect detector is better at catching visually unusual outliers than at catching confidently normal-looking images that are either mislabeled or simply repeated.

\subsection*{Before and after cleaning}
\outputfile{outputs/compare.txt}

\picturefile{figures/clean_curves.png}

\begin{center}\picturefile[0.5\linewidth]{figures/clean_confusion.png}\end{center}

The data for the clean model differs from the data for the first run in two
ways: the cleaning, and the 28 back\_lat\_spread images added after it from a
second download. So this comparison shows the effect of both together, not of
cleaning alone.

\fillin{The test accuracy and the accuracy on your new images, before and
after. If a number went down, why you think it did.}

% ---------------------------------------------------------------------------
\section*{Problem 5. Images from a new source}

\fillin{Where your new images came from, and why you are sure none of them are
among your downloads.}

\outputfile{outputs/overlap.txt}

\fillin{The accuracy on your downloaded test set and on your new images, side
by side.}

The clean model's most confident mistakes on the downloaded test set
(\texttt{results/clean\_worst.png}):

\begin{center}\picturefile{figures/clean_worst.png}\end{center}

\fillin{Five mistakes it was confident about, with the pictures. For each one,
what in the image you think pushed it the wrong way.}

\fillin{What is different about your new images: background, lighting, angle,
distance, what else is in the frame.}

% ---------------------------------------------------------------------------
\section*{Problem 6. The demo}

Address: \url{https://najrgrowth.github.io/mas1004-assignment1/}

\sidebyside{figures/demo_front_double_biceps.png}
           {A front double biceps pose, which is none of my categories. The
            demo page said back\_double\_biceps at 99.1\%.}
           {figures/demo_vacuum_pose_cropped.png}
           {A vacuum pose, which is none of my categories. The demo page said
            back\_double\_biceps at 57.1\%.}

The percentages above were shown by the demo page in the browser, not printed
by code in this repository.

\fillin{What happens when you show it something that is none of your
categories.}

% ---------------------------------------------------------------------------
\section*{Problem 7. What the model looks at}

% Three pictures next to each other, each at most 6 cm high, each with a line
% of text under it: original, changed, control.
\newcommand{\threeone}[2]{%
  \begin{minipage}[t]{0.32\linewidth}\centering
    \IfFileExists{#1}{\includegraphics[width=\linewidth,height=6cm,keepaspectratio]{#1}}%
      {\picturefile{#1}}\par\small #2
  \end{minipage}}
\newcommand{\threeside}[6]{%
  \par\noindent\threeone{#1}{#2}\hfill\threeone{#3}{#4}\hfill\threeone{#5}{#6}\par\medskip}

\subsection*{Guess 1}

\textbf{Guess:} the model chooses back\_double\_biceps based on both arms being up and bent, this is because this is the thing most differentiating from all of the other poses, and also the most recognizable of the back double biceps pose.

\textbf{Prediction:} when the arms are blocked off, the model is less certain about it being back\_double\_biceps.

\textbf{Why:} I chose to cover the arms because this is the most recognizable and differentiating thing about the back double biceps pose. Also because I once put a picture of a front double biceps pose into it and it was 99.1\% certain that it was a back double biceps pose (this came from the demo page, not the code output).

My requests to the agent, word for word, in order:
\begin{quote}
\begin{Verbatim}[breaklines, fontsize=\small]
For each image in data/new_images/back_double_biceps, make two copies:
1. One with a plain grey rectangle covering both of the person's arms.
   Save in data/changed/bdb_arms_covered.
2. One with a plain grey rectangle of the same size covering a part that is
   not the arms (for example the lower back). Save in
   data/changed/bdb_control_covered.
Do not change the originals. Show me one example of each so I can check where
the blocks are.
\end{Verbatim}
\end{quote}
\begin{quote}
\begin{Verbatim}[breaklines, fontsize=\small]
Run the original, arms-covered and control versions of all 7 back double biceps
images through the clean model, prepared the same way as in training. For each
image print the probability for all four classes. Save the output in
results/problem7_guess1.txt and copy it to long_report/outputs/.
\end{Verbatim}
\end{quote}
\begin{quote}
\begin{Verbatim}[breaklines, fontsize=\small]
Please add to src/problem7_predict.py a line per folder that prints the mean
probability of each class over all images, then rerun it for the three back
double biceps folders and save the output again in results/problem7_guess1.txt
and long_report/outputs/.
\end{Verbatim}
\end{quote}

The model's answers below are from \texttt{results/problem7\_guess1.txt}.

\threeside{figures/p7_original_IMG_2306.png}{original: back\_double\_biceps 99.1\%}
           {figures/p7_g1_arms_IMG_2306.png}{arms covered: back\_double\_biceps 71.7\%}
           {figures/p7_g1_control_IMG_2306.png}{control: back\_double\_biceps 95.5\%}

\threeside{figures/p7_original_IMG_9293.jpeg}{original: back\_double\_biceps 95.5\%}
           {figures/p7_g1_arms_IMG_9293.jpeg}{arms covered: side\_triceps 54.7\% (back\_double\_biceps 5.4\%)}
           {figures/p7_g1_control_IMG_9293.jpeg}{control: back\_double\_biceps 97.8\%}

\threeside{figures/p7_original_IMG_0512.jpeg}{original: back\_double\_biceps 98.8\%}
           {figures/p7_g1_arms_IMG_0512.jpeg}{arms covered: side\_triceps 46.3\% (back\_double\_biceps 41.2\%)}
           {figures/p7_g1_control_IMG_0512.jpeg}{control: back\_double\_biceps 94.0\%}

\outputfile{outputs/problem7_guess1.txt}

\textbf{Results:} my prediction was pretty accurate. For the original photos the model was on average 98.1\% certain of it being back double biceps. For the photos with the arms covered, the model was on average 48.1\% certain, with 2 of the images being mistaken for side triceps. The spreading is pretty big, the lowest certainty was 5.4\% and the highest 71.7\%. And for the control covered group, the model was on average 95.7\% certain. (I asked Claude Code to calculate these average percentages.)

\textbf{What I think the model uses (after 1 guess):} I think the arm position is one of the model's most important points for back\_double\_biceps. Covering both arms dropped its average certainty from 98.1\% to 48.1\%, while covering the lower back almost didn't change it (95.7\%). Also, the demo page was 99.1\% sure that a front double biceps photo was a back double biceps. The arm position and silhouette are very similar, and only the direction you face differs.

\textbf{How sure I am:} pretty sure the arms matter, but I am not 100\% positive that it is the most important clue. With the arms covered still 5 of the 7 images were classified as back\_double\_biceps.

\textbf{Next guess:} besides the blocks on the arms I also think the machine looks at the upper back and V shape of the body, so my next guess is that if you just cover the upper back and shoulders, and not the arms, the model will also be less certain of it being a back double biceps.

\subsection*{Guess 2}

\textbf{Guess:} The model chooses back\_double\_biceps based on the back (the upper back, the shoulder blades and the sides of the back, which make the V-shape). If the back is covered while the arms stay visible, the model will be less certain that it is a back\_double\_biceps pose.

\textbf{Prediction:} When the upper part of the back is covered, the model will be less certain about it being back\_double\_biceps, I do think the model will still be a bit more certain than it was with the arms being blocked off, this is because the model will probably still see the back as a certain shape going together with the back double biceps pose.

\textbf{Why:} I chose to cover this part of the body because the V-shape is a very recognizable part of the back double biceps pose and it is also not seen in other poses besides the back lat spread. I also chose it because the model was still 5/7 times correct when covering just the arms.

\textbf{What went wrong and my second request:} The first example was off-centre and covered too little, and the control block was too close to the lower back. So I asked for a redo.

My requests to the agent, word for word, in order:
\begin{quote}
\begin{Verbatim}[breaklines, fontsize=\small]
For each image in data/new_images/back_double_biceps, make two copies:

1. One with a plain grey block covering the upper back and shoulders (the area
between and around the shoulder blades, up to the base of the neck). The block
must NOT overlap the arm areas I covered in guess 1, so the arms stay visible.
Save in data/changed/bdb_upperback_covered.

2. One control copy with a plain grey block of exactly the same size, placed on
a part that is neither the upper back, the shoulders nor the arms (for example
the waist and hips). Save in data/changed/bdb_upperback_control.

Place both blocks so that they are as much as possible inside the centre square
that the model sees, and tell me for each image how much of each block is inside
it. Do not change the originals. Show me one example of each before saving the
rest, so I can check where the blocks are.

Then run src/problem7_predict.py on these folders: the originals,
data/changed/bdb_arms_covered, data/changed/bdb_upperback_covered and
data/changed/bdb_upperback_control. Print the probability for all four classes
per image and the mean per folder. Save the output in
results/problem7_guess2.txt and copy it to long_report/outputs/.
\end{Verbatim}
\end{quote}
\begin{quote}
\begin{Verbatim}[breaklines, fontsize=\small]
The guess 2 example is not good. The upper back block is off-centre (it sits to
the left of the spine) and covers too little, and the control block sits too
close to the lower back. Please redo both:

1. Upper back block: centred on the middle of the body (below the head),
covering the whole back from the base of the neck down to the waist, and wide
enough to go from one armpit to the other, so the shoulder blades and the sides
of the back (the V-shape) are covered. Keep the raised arms visible (upper
arms, elbows, fists). Save in data/changed/bdb_upperback_covered.

2. Control block: exactly the same size, on a part of the image that is clearly
not the back and not the arms, for example the background beside the waist and
hips, or the trousers and legs, and not touching the lower back. Choose the
position per image and tell me where it is. Save in
data/changed/bdb_upperback_control.

Do not save anything yet. First show me one preview sheet with outlines of both
blocks on all 7 images (not just one example), and for each block how much of
it is inside the centre square that the model sees.
\end{Verbatim}
\end{quote}
\begin{quote}
\begin{Verbatim}[breaklines, fontsize=\small]
The preview looks right. For IMG_1481, go with option 2: use it for the back
block test, but leave its control out of the control means, and show that
clearly in the output. Save both folders for all 7 images, then run
problem7_predict.py on the four folders (originals, arms covered, back covered,
back control) into results/problem7_guess2.txt. Print the means per folder over
all 7 images, and also over the 6 images that have a usable control
(without IMG_1481). Copy the output to long_report/outputs/.
\end{Verbatim}
\end{quote}
\begin{quote}
\begin{Verbatim}[breaklines, fontsize=\small]
Please print, for each image, the area in pixels of the two arm blocks together
(guess 1) and of the back block (guess 2), and the ratio between them. Save it in
results/problem7_block_sizes.txt and copy it to long_report/outputs/.
\end{Verbatim}
\end{quote}

The model's answers below are from \texttt{results/problem7\_guess2.txt}.

\threeside{figures/p7_original_IMG_2306.png}{original: back\_double\_biceps 99.1\%}
           {figures/p7_g2_back_IMG_2306.png}{back covered: back\_double\_biceps 93.3\%}
           {figures/p7_g2_control_IMG_2306.png}{control: back\_double\_biceps 98.9\%}

\threeside{figures/p7_original_IMG_9293.jpeg}{original: back\_double\_biceps 95.5\%}
           {figures/p7_g2_back_IMG_9293.jpeg}{back covered: back\_double\_biceps 79.4\%}
           {figures/p7_g2_control_IMG_9293.jpeg}{control: back\_double\_biceps 97.3\%}

\threeside{figures/p7_original_IMG_0512.jpeg}{original: back\_double\_biceps 98.8\%}
           {figures/p7_g2_back_IMG_0512.jpeg}{back covered: back\_double\_biceps 67.9\%}
           {figures/p7_g2_control_IMG_0512.jpeg}{control: back\_double\_biceps 98.5\%}

\outputfile{outputs/problem7_guess2.txt}
\outputfile{outputs/problem7_block_sizes.txt}

\textbf{Results:} My prediction was partly right. The certainty dropped, from 98.1\% to 83.9\% (84.1\% without IMG\_1481) and the drop was smaller than with the arms being covered (48.1\%). But I expected the model to only be a bit more certain than with the arms covered, and the difference is actually large. None of the 7 images changed class and the values ranged from 67.9\% to 93.3\%. The control stayed at 98.1\%. The back block covered less area than the two arm blocks together on every image (between 0.39 and 0.74 times as much, see \texttt{results/problem7\_block\_sizes.txt}). Still, the much smaller drop suggests that the arms matter more than the back.

\textbf{What I think the model uses:} I still believe the model focuses mostly on the arms, but the upper back still plays a role too.

\textbf{How sure I am:} I am fairly sure that covering the arms had a bigger effect than covering the back for these 7 images. I cannot say that the arms matter more by themselves, because the blocks were not equally large. I also did not cover the arms and back together, and I tested only 7 images and one model. The control for IMG\_1481 was left out (16\% of its block was inside the model's view), the control for IMG\_2359 was only 67\% inside it, and for three images the control block lay on a dark background.

\subsection*{Guess 3}

\textbf{My guess:} the bigger drop with the arms covered was caused by the arms themselves, not by the bigger covered area. If I cover a tall block in the middle of the body, from the top of the head to the waist, with the same area as the two arm blocks together the drop will be smaller.

\textbf{Prediction:} the average certainty for back\_double\_biceps will stay above the 48.1\% from the arms test, and fewer than 2 of the 7 images will change class.

\textbf{Why:} In guess 1 and 2 the blocks differed in size, so I could not tell if the arms or the area caused the drop. With the same area as the arm blocks, that doubt is gone. I put it in the middle of the body, where the arms are not.

The block was too wide and it would overlap with the shoulders and a part of the arms so extra prompt.

This also did not really work, as now on 5 of 7 pictures there was no space left for the control block, and on 2 of those the covered block ran off the bottom of the photo. Therefore I sent this prompt.

My requests to the agent, word for word, in order:
\begin{quote}
\begin{Verbatim}[breaklines, fontsize=\small]
For each image in data/new_images/back_double_biceps, make two copies:

1. One with a single plain grey rectangle whose area is equal to the two guess 1
arm blocks together (use the pixel areas in src/problem7_block_sizes.py, within
1%). Make it tall: centred on the spine, from the top of the head downwards
over the neck and back. Save in data/changed/bdb_samearea_covered.

2. One control copy with a grey rectangle of exactly the same size and shape,
on a part of the image that is clearly not the body (background or trousers).
Choose the position per image and tell me where it is. Save in
data/changed/bdb_samearea_control.

Keep both blocks as much as possible inside the centre square that the model
sees, and tell me per image how much of each block is inside it. Do not change
the originals. Do not save anything yet: first show me one preview sheet with
outlines of both blocks on all 7 images. If a control block does not fit on an
image, tell me instead of moving it somewhere else.

Then add the new block positions to a script in src/ (so the areas come from
code in git), run src/problem7_predict.py on the originals, arms covered,
same-area covered and same-area control, and print per-image probabilities and
the mean per folder into results/problem7_guess3.txt. Copy it to
long_report/outputs/.
\end{Verbatim}
\end{quote}
\begin{quote}
\begin{Verbatim}[breaklines, fontsize=\small]
Make the same-area block narrow enough so it does NOT overlap the guess 1 arm
blocks (keep the area equal to the two arm blocks, within 1%). It will be longer
and may go onto the trousers, that is fine. Keep it centred on the spine, starting
at the top of the head. Re-plan the control blocks for this new shape. For every
image print how much of the covered block and of the control block is inside
the centre square the model sees. Do not save anything yet: show me a new
preview of all 7 images first. If the covered block is mostly outside the
model's view for several images, tell me instead of changing it.
\end{Verbatim}
\end{quote}
\begin{quote}
\begin{Verbatim}[breaklines, fontsize=\small]
Go back to the first plan: same area as the two arm blocks (within 1%), centred
on the spine, from the top of the head down to the waist, so it may overlap the
guess 1 arm blocks. Use the control positions from that first plan. Leave the
controls of IMG_1481 and IMG_2359 out of the control means (print their answers,
marked), as in guess 2. Put the block positions in a script in src/ and make it
print per image the area, how much of each block is inside the model's view, and
the overlap with the arm blocks in pixels and as a share of the arm area. Show
me the final preview of all 7 images before saving anything. Then save both
folders, run problem7_predict.py on the originals, arms covered, same-area
covered and same-area control, and save the output in
results/problem7_guess3.txt and long_report/outputs/.
\end{Verbatim}
\end{quote}
\begin{quote}
\begin{Verbatim}[breaklines, fontsize=\small]
Yes, go ahead. Save the two folders with --write and check the originals are
unchanged. Run problem7_predict.py on the originals, arms covered, same-area
covered and same-area control. Mark the controls of IMG_1481 and IMG_2359 and
leave them out of the control means (print their answers anyway). Save the
output in results/problem7_guess3.txt and copy it to long_report/outputs/.
Then show me the means per folder and how many of the 7 images change class.
\end{Verbatim}
\end{quote}
\begin{quote}
\begin{Verbatim}[breaklines, fontsize=\small]
Yes. Add a line per folder to problem7_predict.py that prints how many images
change class, and rerun guess 3 into results/problem7_guess3.txt (and copy it
to long_report/outputs/). Save the table from problem7_samearea.py to
results/problem7_samearea.txt and copy it to long_report/outputs/. Then commit
src/problem7_samearea.py and these outputs and push.
\end{Verbatim}
\end{quote}

The model's answers below are from \texttt{results/problem7\_guess3.txt}.

\threeside{figures/p7_original_IMG_2306.png}{original: back\_double\_biceps 99.1\%}
           {figures/p7_g3_samearea_IMG_2306.png}{same area covered: back\_double\_biceps 87.3\%}
           {figures/p7_g3_control_IMG_2306.png}{control: back\_double\_biceps 99.0\%}

\threeside{figures/p7_original_IMG_9293.jpeg}{original: back\_double\_biceps 95.5\%}
           {figures/p7_g3_samearea_IMG_9293.jpeg}{same area covered: back\_lat\_spread 43.9\% (back\_double\_biceps 38.3\%)}
           {figures/p7_g3_control_IMG_9293.jpeg}{control: back\_double\_biceps 98.2\%}

\threeside{figures/p7_original_IMG_0512.jpeg}{original: back\_double\_biceps 98.8\%}
           {figures/p7_g3_samearea_IMG_0512.jpeg}{same area covered: back\_double\_biceps 66.3\%}
           {figures/p7_g3_control_IMG_0512.jpeg}{control: back\_double\_biceps 98.4\%}

\outputfile{outputs/problem7_samearea.txt}
\outputfile{outputs/problem7_guess3.txt}

\textbf{Results:} covering the same area in the middle of the body gives a certainty of 71.1\%, while covering the arms gives 48.1\%. The grey block in the last round also covered 5.7\% to 31.1\% of the arm/shoulder area, this should work against my guess but still the difference stays big.

\textbf{What I think the model uses:} my prediction came true, 71.1\% compared to 48.1\%, and 1 of 7 images changed class instead of 2. This suggests that the arms count for more than their area alone. But the middle block still hurts the model, which means it is not only the arms.

\textbf{How sure I am:} I am fairly sure that the arms matter more than the area alone, because the gap in score is large and the control stays similar. I am not fully sure how much the rest of the body matters, but I did cover the arms, upper back and the middle of the body. Yet I would not claim exact numbers.

\textbf{Limits:} the control group did only exist out of 5 images, this was because IMG\_1481 was only 9\% visible and IMG\_2359 35\%. Also in IMG\_9293, 16\% of the covered block lay outside the square.


\subsection*{What I now think the model uses}
What I now think the model uses after all 3 guesses, is that the arms are most likely the most important part. Covering the arms lowered the mean certainty for back\_double\_biceps the most (98.1\% to 48.1\%). Covering the same area in the middle of the body lowered it less (71.1\%), and covering the upper back, with a block about half as big, lowered it the least (83.9\%). The controls hardly changed anything (95.7\%, 98.1\% and 98.4\%), so a grey block on its own does not confuse the model. I think the model mostly uses the raised, bent arms and how they sit next to the body. This fits my guess from before the training. The rest of the body also matters, because the same-area block in the middle still cost quite some points, but less than the arms. The upper back, the V-shape, seems to play a smaller role than I thought at first.

I am fairly sure that the arms matter more than the area alone, because the gap is large and the controls stay similar. I am less sure how much the back matters by itself, because the back block was smaller and the middle block overlapped the arm blocks in all 7 images. I tested 7 images and one model, I did not cover several body parts at once, and the blocks were placed by eye, so I would not claim exact numbers.

% ---------------------------------------------------------------------------
\section*{Appendix. Output of check.py}
\outputfile{outputs/check.txt}

\end{document}
